% ========================================================================
%  MinSTEP Kernel Abstraction Layer
%  Reference Manual, Version 1.0
%
%  Typeset in TeX to evoke the NeXT Computer, Inc. technical
%  reference manuals of the late 1980s.
% ========================================================================

\documentclass[11pt,twoside]{article}

\usepackage[margin=1in]{geometry}
\usepackage{courier}
\usepackage{url}
\usepackage{tabularx}
\usepackage{longtable}
\usepackage{fancyhdr}

% -----------------------------------------------------------------------
%  Page style
% -----------------------------------------------------------------------

\pagestyle{fancy}
\fancyhf{}
\fancyhead[LE]{\scshape\small MINSTEP Kernel Abstraction Layer}
\fancyhead[RO]{\scshape\small Reference Manual}
\fancyfoot[CE]{\small\thepage}
\fancyfoot[CO]{\small\thepage}
\renewcommand{\headrulewidth}{0.4pt}
\renewcommand{\footrulewidth}{0pt}

% -----------------------------------------------------------------------
%  Commands
% -----------------------------------------------------------------------

\newcommand{\KAL}{KAL}
\newcommand{\Mango}{Mango}
\newcommand{\MinSTEP}{MinSTEP}
\newcommand{\NeXTSTEP}{NeXTSTEP}
\newcommand{\code}[1]{{\tt #1}}
\newcommand{\meta}[1]{{\itshape #1}}
\newcommand{\ie}{{\itshape i.e.}}

% ========================================================================

\begin{document}

% -----------------------------------------------------------------------
%  Title page
% -----------------------------------------------------------------------

\thispagestyle{empty}
\begin{center}
\vspace*{2in}

{\LARGE\bfseries MinSTEP}\\[0.3in]
{\large\itshape An Open-Source NeXTSTEP 0.8-like Operating Environment}\\[0.5in]
{\LARGE\bfseries Kernel Abstraction Layer}\\[0.3in]
{\Large\bfseries Reference Manual}\\[0.8in]
{\large Version 1.0}\\[0.5in]
\vfill
{\large Miguel V.\ Mesquita}\\[0.2in]
{\large July 2026}

\end{center}

\newpage

% -----------------------------------------------------------------------
%  Copyright page
% -----------------------------------------------------------------------

\thispagestyle{empty}
\vspace*{8in}

\begin{center}
\small
Copyright {\copyright} 2026 Miguel V.\ Mesquita.\\
All rights reserved.

\medskip

This document is part of the MinSTEP project and is distributed
under the BSD License.

\medskip

MinSTEP is an open-source project inspired by NeXTSTEP 0.8--2.2.\\
It provides a usermode operating environment with its own kernel
and Mach IPC.

\medskip

NeXTSTEP is a trademark of NeXT, Inc.\\
Mach is a trademark of Carnegie Mellon University.

\bigskip

{\bf MinSTEP}\\
Licensed under the BSD License.
\end{center}

\newpage

% -----------------------------------------------------------------------
%  Table of contents
% -----------------------------------------------------------------------

\tableofcontents
\newpage

% ========================================================================
\section{Introduction}
% ========================================================================

This document describes the MinSTEP Kernel Abstraction Layer (\KAL),
a thin, portable abstraction over host operating system primitives.
The \KAL\ allows the \Mango\ nanokernel to run on any supported
platform without modifying kernel subsystem code.

% ========================================================================
\subsection{Overview}
% ========================================================================

The \Mango\ nanokernel is a usermode process that emulates Mach kernel
abstractions---ports, tasks, IPC, and binary loading---on top of a
host operating system.  In its original form, the kernel called
POSIX APIs directly (\code{fork}, \code{socketpair}, \code{poll},
\code{pthread\_mutex}, \code{tcgetattr}, \code{clock\_gettime}, and
so on).  This tied the kernel to Unix-like systems.

The \KAL\ replaces these direct calls with portable function
interfaces.  Each interface has one or more backend implementations,
selected at compile time by platform detection macros.  The current
release includes a POSIX backend for Linux, macOS, FreeBSD, and other
Unix-like systems.  Additional backends (\ie\ Win32) may be added
in the future.

The translation model is:

\begin{center}
{\tt kernel subsystem}~$\longrightarrow$~{\bfseries [KAL API]}~$\longrightarrow$~{\bfseries [platform backend]}~$\longrightarrow$~{\tt host OS}
\end{center}

The \KAL\ is written in standard C (C99) and compiles on any
platform that provides a C compiler.  It consists of 10 header files
and 9 implementation files (for the POSIX backend), totalling
approximately 1{,}200 lines of code.

% ========================================================================
\subsection{Design Principles}
% ========================================================================

The \KAL\ follows four design principles:

\begin{enumerate}
  \item {\bf Thin wrappers.}  Every \KAL\ function does little more
        than call the corresponding host API.  There is no hidden
        state, no middleware, and no performance overhead beyond one
        function call.

  \item {\bf Compile-time selection.}  Platform backends are selected
        by {\tt \#ifdef} in {\tt kal\_platform.h}.  There are no
        function pointer tables or runtime dispatch.  The cost is
        zero at run time.

  \item {\bf Graceful fallback.}  The Compat headers ({\tt lock.h},
        {\tt sched\_prim.h}, {\tt zalloc.h}) detect whether the \KAL\
        is available and fall back to direct POSIX calls if it is not.
        This means the kernel can still build without the \KAL\ on
        supported platforms.

  \item {\bf Minimal surface area.}  The \KAL\ abstracts only the
        host APIs that the kernel actually uses.  It does not attempt
        to be a general-purpose OS abstraction layer.
\end{enumerate}

% ========================================================================
\subsection{Supported Platforms}
% ========================================================================

The following platforms are detected automatically by
{\tt kal\_platform.h}:

\medskip
\noindent
\begin{tabularx}{\textwidth}{lX}
{\bf Platform}       & {\bf Macro} \\[2pt]
Linux                & {\tt KAL\_PLATFORM\_LINUX} \\
macOS (Darwin)       & {\tt KAL\_PLATFORM\_DARWIN} \\
FreeBSD              & {\tt KAL\_PLATFORM\_FREEBSD} \\
NetBSD               & {\tt KAL\_PLATFORM\_NETBSD} \\
OpenBSD              & {\tt KAL\_PLATFORM\_OPENBSD} \\
Other Unix           & {\tt KAL\_PLATFORM\_POSIX} \\
Windows              & {\tt KAL\_PLATFORM\_WINDOWS} \\
\end{tabularx}

All Unix-like platforms define {\tt KAL\_PLATFORM\_POSIX} and share
the POSIX backend.  Windows detection is present but the Win32
backend is not yet implemented.

% ========================================================================
\newpage
\section{Header Files}
% ========================================================================

The \KAL\ is accessed through a single master header:

\begin{quote}
{\tt \#include <kal/kal.h>}
\end{quote}

This pulls in all sub-headers.  Individual headers may also be
included directly if only a subset of the \KAL\ is needed.

\bigskip
\noindent
\begin{tabularx}{\textwidth}{lX}
{\bf Header} & {\bf Purpose} \\[2pt]
{\tt kal.h}            & Master include---pulls in everything. \\
{\tt kal\_platform.h}  & Platform detection and feature flags. \\
{\tt kal\_process.h}   & Process creation, termination, environment. \\
{\tt kal\_ipc.h}       & Socket pairs, read/write, poll. \\
{\tt kal\_memory.h}    & Memory allocation. \\
{\tt kal\_signal.h}    & Signal handler registration. \\
{\tt kal\_thread.h}    & Mutexes, condition variables, yield. \\
{\tt kal\_terminal.h}  & Terminal save/restore/raw mode. \\
{\tt kal\_filesys.h}   & File I/O, temporary files, directories. \\
{\tt kal\_system.h}    & System information, memory, machine type. \\
{\tt kal\_time.h}      & Monotonic clock and timestamps. \\
\end{tabularx}

% ========================================================================
\newpage
\section{Process Management}
\label{sec:process}
% ========================================================================

The process interface abstracts process creation, termination, and
environment variable access.

% ========================================================================
\subsection{Types}
% ========================================================================

\begin{quote}
\begin{verbatim}
typedef int kal_pid_t;
\end{verbatim}
\end{quote}

A process identifier.  On POSIX systems this is {\tt pid\_t}.

% ========================================================================
\subsection{Process Identity}
% ========================================================================

\begin{quote}
\begin{verbatim}
kal_pid_t kal_getpid(void);
\end{verbatim}
\end{quote}

Returns the PID of the calling process.

% ========================================================================
\subsection{Process Creation}
% ========================================================================

\begin{quote}
\begin{verbatim}
typedef struct kal_spawn_args {
    const char     *exec_path;
    const char    **argv;
    const char    **envp;
} kal_spawn_args_t;

int kal_process_spawn(kal_spawn_args_t *args,
                      kal_pid_t *out_pid);
\end{verbatim}
\end{quote}

Spawns a new process.  The child executes {\tt exec\_path} with the
given argument vector.  If {\tt envp} is {\tt NULL}, the child
inherits the parent's environment.

On success, {\tt *out\_pid} is set to the child's PID and the
function returns $0$.  On failure it returns $-1$.

\medskip
\noindent
{\bf Example:}
\begin{quote}
\begin{verbatim}
kal_pid_t pid;
const char *argv[] = { "/sbin/init", NULL };

kal_spawn_args_t args = {
    .exec_path = "/sbin/init",
    .argv = argv,
    .envp = NULL
};

if (kal_process_spawn(&args, &pid) == 0) {
    klog_info("init launched (pid %d)\n", pid);
}
\end{verbatim}
\end{quote}

% ========================================================================
\subsection{Process Termination}
% ========================================================================

\begin{quote}
\begin{verbatim}
int kal_process_kill(kal_pid_t pid, int sig);
\end{verbatim}
\end{quote}

Sends signal {\tt sig} to process {\tt pid}.  Returns $0$ on
success.

\begin{quote}
\begin{verbatim}
void kal_process_exit(int code)
    __attribute__((noreturn));
\end{verbatim}
\end{quote}

Terminates the calling process with exit code {\tt code}.

% ========================================================================
\subsection{Process Reaping}
% ========================================================================

\begin{quote}
\begin{verbatim}
kal_pid_t kal_process_reap(int *out_status);
\end{verbatim}
\end{quote}

Reaps a terminated child process without blocking ({\tt WNOHANG}).
Returns the child PID, or $0$ if no children have exited.  If
{\tt out\_status} is not {\tt NULL}, the exit status is stored there.

\begin{quote}
\begin{verbatim}
kal_pid_t kal_process_wait(kal_pid_t pid,
                           int *out_status);
\end{verbatim}
\end{quote}

Waits for a specific child process to exit (blocking).  Returns the
child PID.

% ========================================================================
\subsection{Environment Variables}
% ========================================================================

\begin{quote}
\begin{verbatim}
const char *kal_env_get(const char *name);
int kal_env_set(const char *name,
                const char *value,
                int overwrite);
\end{verbatim}
\end{quote}

{\tt kal\_env\_get} returns the value of environment variable
{\tt name}, or {\tt NULL} if it is not set.

{\tt kal\_env\_set} sets or updates an environment variable.  If
{\tt overwrite} is nonzero, an existing value is replaced.  Returns
$0$ on success.

% ========================================================================
\subsection{File Executability Check}
% ========================================================================

\begin{quote}
\begin{verbatim}
int kal_can_exec(const char *path);
\end{verbatim}
\end{quote}

Returns $1$ if {\tt path} exists and is executable by the current
user, $0$ otherwise.

% ========================================================================
\newpage
\section{IPC and Socket Primitives}
\label{sec:ipc}
% ========================================================================

The IPC interface abstracts socket pairs and poll-based I/O, which
the \Mango\ kernel uses to emulate Mach port message passing.

% ========================================================================
\subsection{Types}
% ========================================================================

\begin{quote}
\begin{verbatim}
typedef int kal_fd_t;

#define KAL_FD_INVALID  (-1)
\end{verbatim}
\end{quote}

A file descriptor handle.  {\tt KAL\_FD\_INVALID} represents an
invalid descriptor.

\begin{quote}
\begin{verbatim}
typedef struct kal_pollfd {
    kal_fd_t    fd;
    uint32_t    events;
    uint32_t    revents;
} kal_pollfd_t;
\end{verbatim}
\end{quote}

A poll descriptor.  The {\tt events} field specifies the requested
event mask; {\tt revents} receives the returned events.

\medskip
\noindent
\begin{tabularx}{\textwidth}{lX}
{\bf Constant} & {\bf Meaning} \\[2pt]
{\tt KAL\_POLLIN}  & Data is available for reading. \\
{\tt KAL\_POLLOUT} & Writing will not block. \\
{\tt KAL\_POLLERR} & Error condition. \\
{\tt KAL\_POLLHUP} & Hang-up. \\
\end{tabularx}

% ========================================================================
\subsection{Socket Pairs}
% ========================================================================

\begin{quote}
\begin{verbatim}
int kal_socket_pair(kal_fd_t *out_fd0,
                    kal_fd_t *out_fd1);
\end{verbatim}
\end{quote}

Creates a connected pair of Unix domain socket file descriptors.
Returns $0$ on success, $-1$ on error.

% ========================================================================
\subsection{I/O}
% ========================================================================

\begin{quote}
\begin{verbatim}
ssize_t kal_write(kal_fd_t fd,
                  const void *buf,
                  size_t count);

ssize_t kal_read(kal_fd_t fd,
                 void *buf,
                 size_t count);
\end{verbatim}
\end{quote}

{\tt kal\_write} writes up to {\tt count} bytes from {\tt buf} to
the file descriptor {\tt fd}.  Returns the number of bytes written,
or $-1$ on error.

{\tt kal\_read} reads up to {\tt count} bytes into {\tt buf}.
Returns the number of bytes read, $0$ on end-of-file, or $-1$ on
error.

\begin{quote}
\begin{verbatim}
int kal_close(kal_fd_t fd);
\end{verbatim}
\end{quote}

Closes a file descriptor.

% ========================================================================
\subsection{Polling}
% ========================================================================

\begin{quote}
\begin{verbatim}
int kal_poll(kal_pollfd_t *fds,
             unsigned int nfds,
             int timeout_ms);
\end{verbatim}
\end{quote}

Waits for events on up to {\tt nfds} file descriptors.  Returns the
number of ready descriptors, $0$ on timeout, or $-1$ on error.  At
most 64 descriptors may be polled per call.

% ========================================================================
\newpage
\section{Memory Management}
\label{sec:memory}
% ========================================================================

The memory interface wraps the host allocator.  On platforms without
a standard C library, these may be backed by a custom heap.

\begin{quote}
\begin{verbatim}
void *kal_malloc(size_t size);
void *kal_calloc(size_t count, size_t size);
void *kal_realloc(void *ptr, size_t new_size);
void  kal_free(void *ptr);
\end{verbatim}
\end{quote}

These follow standard C semantics.  {\tt kal\_calloc} returns
zero-initialized memory.  {\tt kal\_free} accepts {\tt NULL} safely.

% ========================================================================
\newpage
\section{Signal Handling}
\label{sec:signal}
% ========================================================================

The signal interface provides portable signal registration.

\begin{quote}
\begin{verbatim}
#define KAL_SIGTERM  15
#define KAL_SIGINT    2
#define KAL_SIGCHLD  17
#define KAL_SIG_IGN  ((kal_sighandler_t)1)

typedef void (*kal_sighandler_t)(int sig);
\end{verbatim}
\end{quote}

The signal number constants are defined to their POSIX values.
On Windows or other platforms, the backend maps them to the
appropriate native codes.

\begin{quote}
\begin{verbatim}
kal_sighandler_t kal_signal(
    int sig,
    kal_sighandler_t handler);
\end{verbatim}
\end{quote}

Installs a signal handler.  Returns the previous handler.

\medskip
\noindent
{\bf Example:}
\begin{quote}
\begin{verbatim}
static volatile int running = 1;

static void on_shutdown(int sig) {
    (void)sig;
    running = 0;
}

kal_signal(KAL_SIGTERM, on_shutdown);
kal_signal(KAL_SIGINT,  on_shutdown);
kal_signal(KAL_SIGCHLD, KAL_SIG_IGN);
\end{verbatim}
\end{quote}

% ========================================================================
\newpage
\section{Threading and Synchronization}
\label{sec:thread}
% ========================================================================

The threading interface provides mutexes, condition variables, and
thread yield.  These are used by the Mach4 compat layer
({\tt lock.h}, {\tt sched\_prim.h}) to implement kernel
synchronization primitives.

% ========================================================================
\subsection{Mutex}
% ========================================================================

\begin{quote}
\begin{verbatim}
typedef struct kal_mutex {
    void *_impl;
} kal_mutex_t;
\end{verbatim}
\end{quote}

An opaque mutex handle.  The {\tt \_impl} pointer is
allocated by {\tt kal\_mutex\_init} and holds the
backend-specific state (on POSIX, a {\tt pthread\_mutex\_t}).

\begin{quote}
\begin{verbatim}
int kal_mutex_init(kal_mutex_t *m);
int kal_mutex_destroy(kal_mutex_t *m);
int kal_mutex_lock(kal_mutex_t *m);
int kal_mutex_trylock(kal_mutex_t *m);
int kal_mutex_unlock(kal_mutex_t *m);
\end{verbatim}
\end{quote}

All functions return $0$ on success and $-1$ on error.
{\tt kal\_mutex\_trylock} returns $0$ if the lock was acquired
and $-1$ if it was not.

% ========================================================================
\subsection{Condition Variable}
% ========================================================================

\begin{quote}
\begin{verbatim}
typedef struct kal_cond {
    void *_impl;
} kal_cond_t;
\end{verbatim}
\end{quote}

\begin{quote}
\begin{verbatim}
int kal_cond_init(kal_cond_t *c);
int kal_cond_destroy(kal_cond_t *c);
int kal_cond_wait(kal_cond_t *c, kal_mutex_t *m);
int kal_cond_signal(kal_cond_t *c);
int kal_cond_broadcast(kal_cond_t *c);
\end{verbatim}
\end{quote}

{\tt kal\_cond\_wait} atomically releases the mutex {\tt m} and
blocks on the condition variable {\tt c}.  When it returns, the
mutex is re-acquired.

{\tt kal\_cond\_signal} wakes one waiter.
{\tt kal\_cond\_broadcast} wakes all waiters.

% ========================================================================
\subsection{Thread Yield}
% ========================================================================

\begin{quote}
\begin{verbatim}
void kal_yield(void);
\end{verbatim}
\end{quote}

Yields the current thread's time slice to the host scheduler.

% ========================================================================
\newpage
\section{Terminal I/O}
\label{sec:terminal}
% ========================================================================

The terminal interface provides raw-mode manipulation for the boot
prompt and kernel console.

\begin{quote}
\begin{verbatim}
typedef struct kal_terminal {
    void *_impl;
} kal_terminal_t;
\end{verbatim}
\end{quote}

An opaque terminal handle that stores the saved terminal state.

\begin{quote}
\begin{verbatim}
int kal_terminal_save(kal_terminal_t *t);
int kal_terminal_restore(kal_terminal_t *t);
int kal_terminal_raw(kal_terminal_t *t);
\end{verbatim}
\end{quote}

{\tt kal\_terminal\_save} captures the current terminal settings.
{\tt kal\_terminal\_restore} restores previously saved settings.
{\tt kal\_terminal\_raw} puts the terminal into non-canonical,
no-echo mode suitable for character-by-character input.

\medskip
\noindent
{\bf Example:}
\begin{quote}
\begin{verbatim}
kal_terminal_t term;
int saved = (kal_terminal_save(&term) == 0);

/* ... run kernel ... */

if (saved)
    kal_terminal_restore(&term);
\end{verbatim}
\end{quote}

% ========================================================================
\newpage
\section{Filesystem Operations}
\label{sec:filesys}
% ========================================================================

The filesystem interface abstracts file I/O, temporary files,
directory creation, and permissions.  These are used primarily by
the {\tt .mach} binary loader.

% ========================================================================
\subsection{File Handle}
% ========================================================================

\begin{quote}
\begin{verbatim}
typedef struct kal_file {
    void *_impl;
} kal_file_t;
\end{verbatim}
\end{quote}

An opaque file handle wrapping the host {\tt FILE*} (on POSIX).

% ========================================================================
\subsection{High-Level File I/O}
% ========================================================================

\begin{quote}
\begin{verbatim}
kal_file_t *kal_fopen(const char *path,
                      const char *mode);
int kal_fclose(kal_file_t *f);

size_t kal_fread(void *buf, size_t size,
                 size_t count, kal_file_t *f);
size_t kal_fwrite(const void *buf, size_t size,
                  size_t count, kal_file_t *f);

int kal_fseek(kal_file_t *f,
              long offset, int whence);
\end{verbatim}
\end{quote}

These follow standard C {\tt FILE} semantics.  The {\tt whence}
argument uses {\tt KAL\_SEEK\_SET}, {\tt KAL\_SEEK\_CUR}, and
{\tt KAL\_SEEK\_END}.

% ========================================================================
\subsection{Low-Level File I/O}
% ========================================================================

\begin{quote}
\begin{verbatim}
int kal_open(const char *path,
             int flags, int mode);
ssize_t kal_fd_write(int fd,
                     const void *buf,
                     size_t count);
int kal_fd_close(int fd);
\end{verbatim}
\end{quote}

These provide direct file descriptor access for temporary file
extraction.

% ========================================================================
\subsection{Temporary Files}
% ========================================================================

\begin{quote}
\begin{verbatim}
int kal_mkstemp(char *template);
\end{verbatim}
\end{quote}

Creates and opens a unique temporary file.  The {\tt template}
string is modified in place (the last six characters are replaced
with unique characters).  Returns the file descriptor.

% ========================================================================
\subsection{Directory and Metadata Operations}
% ========================================================================

\begin{quote}
\begin{verbatim}
int kal_mkdir(const char *path, int mode);
int kal_chmod(const char *path, int mode);
int kal_unlink(const char *path);
int kal_access(const char *path, int mode);
\end{verbatim}
\end{quote}

These provide portable directory creation, permission changes, file
deletion, and existence checks.

% ========================================================================
\newpage
\section{System Information}
\label{sec:system}
% ========================================================================

The system interface provides queries for hardware and OS properties.

\begin{quote}
\begin{verbatim}
size_t   kal_get_page_size(void);
uint64_t kal_get_total_memory(void);
uint64_t kal_get_available_memory(void);
\end{verbatim}
\end{quote}

{\tt kal\_get\_page\_size} returns the system page size in bytes.
{\tt kal\_get\_total\_memory} and {\tt kal\_get\_available\_memory}
return physical memory sizes in bytes, or $0$ if unknown.

\begin{quote}
\begin{verbatim}
const char *kal_get_machine_type(void);
const char *kal_get_os_name(void);
\end{verbatim}
\end{quote}

{\tt kal\_get\_machine\_type} returns a string identifying the
architecture (e.g.\ {\tt "x86\_64"}, {\tt "aarch64"}).
{\tt kal\_get\_os\_name} returns the OS name and version string.

Both return pointers to static buffers that are initialized on first
call.

\begin{quote}
\begin{verbatim}
void kal_usleep(unsigned int usec);
\end{verbatim}
\end{quote}

Sleeps for the given number of microseconds.

% ========================================================================
\newpage
\section{Timing and Clock}
\label{sec:time}
% ========================================================================

The timing interface provides a monotonic clock for kernel
timestamps.

\begin{quote}
\begin{verbatim}
typedef struct kal_timespec {
    uint64_t    tv_sec;
    uint32_t    tv_nsec;
} kal_timespec_t;
\end{verbatim}
\end{quote}

\begin{quote}
\begin{verbatim}
uint64_t kal_clock_monotonic_ns(void);
void     kal_clock_monotonic(kal_timespec_t *ts);
double   kal_clock_seconds(void);
\end{verbatim}
\end{quote}

{\tt kal\_clock\_monotonic\_ns} returns the time in nanoseconds
since an arbitrary epoch (typically boot time).  The epoch is fixed
for the lifetime of the process.

{\tt kal\_clock\_monotonic} fills a {\tt kal\_timespec\_t} with
the same value split into seconds and nanoseconds.

{\tt kal\_clock\_seconds} returns the time as a double-precision
floating-point value in seconds.

% ========================================================================
\newpage
\section{Platform Configuration}
\label{sec:platform}
% ========================================================================

% ========================================================================
\subsection{Platform Detection}
% ========================================================================

The {\tt kal\_platform.h} header detects the host operating system
and defines exactly one {\tt KAL\_PLATFORM\_*} macro.

% ========================================================================
\subsection{Feature Flags}
% ========================================================================

The following boolean macros indicate which host features are
available:

\bigskip
\noindent
\begin{tabularx}{\textwidth}{lX}
{\bf Macro} & {\bf Meaning} \\[2pt]
{\tt KAL\_HAS\_FORK}       & {\tt fork()} is available. \\
{\tt KAL\_HAS\_EXEC}       & {\tt exec} family is available. \\
{\tt KAL\_HAS\_SOCKETPAIR} & {\tt socketpair()} is available. \\
{\tt KAL\_HAS\_POLL}       & {\tt poll()} is available. \\
{\tt KAL\_HAS\_PTHREAD}    & POSIX threads are available. \\
{\tt KAL\_HAS\_TERMIOS}    & Terminal I/O is available. \\
{\tt KAL\_HAS\_CLOCK\_MONO} & Monotonic clock is available. \\
{\tt KAL\_HAS\_MKSTEMP}    & {\tt mkstemp()} is available. \\
{\tt KAL\_HAS\_ENVIRON}    & Environment variables are available. \\
\end{tabularx}

The POSIX backend defines all flags to $1$.  The Windows backend
defines them all to $0$ except {\tt KAL\_HAS\_ENVIRON}.

% ========================================================================
\subsection{Version}
% ========================================================================

\begin{quote}
\begin{verbatim}
#define KAL_VERSION_MAJOR   1
#define KAL_VERSION_MINOR   0
#define KAL_VERSION_PATCH   0
#define KAL_VERSION_STRING  "1.0.0"
\end{verbatim}
\end{quote}

% ========================================================================
\newpage
\section{Porting the \KAL\ to a New Platform}
\label{sec:porting}
% ========================================================================

To add support for a new operating system:

\begin{enumerate}
  \item {\bf Add platform detection.}
        In {\tt kal\_platform.h}, add a new block:
\begin{quote}
\begin{verbatim}
#elif defined(__myplatform__)
    #define KAL_PLATFORM_MYOS    1
    #define KAL_PLATFORM_NAME    "myos"
\end{verbatim}
\end{quote}

  \item {\bf Set feature flags.}
        Define {\tt KAL\_HAS\_*} macros for each feature the
        platform supports:
\begin{quote}
\begin{verbatim}
    #define KAL_HAS_FORK         1
    #define KAL_HAS_EXEC         1
    #define KAL_HAS_SOCKETPAIR   0
    #define KAL_HAS_POLL         1
    #define KAL_HAS_PTHREAD      0
    #define KAL_HAS_TERMIOS      0
    #define KAL_HAS_CLOCK_MONO   1
    #define KAL_HAS_MKSTEMP      1
    #define KAL_HAS_ENVIRON      1
\end{verbatim}
\end{quote}

  \item {\bf Create backend source files.}
        Create a new directory, e.g.\ {\tt kal/myos/}, and implement
        one {\tt .c} file per \KAL\ subsystem:
        {\tt kal\_myos\_process.c},
        {\tt kal\_myos\_ipc.c},
        {\tt kal\_myos\_memory.c}, etc.

        Each file must implement every function declared in the
        corresponding {\tt kal\_*\_h} header.

  \item {\bf Update the Makefile.}
        Add the new source files to the {\tt KAL\_SRC} and
        {\tt KAL\_OBJ} variables in the kernel {\tt Makefile}.

  \item {\bf Test.}
        Build the kernel and verify that all subsystems initialize
        correctly.
\end{enumerate}

% ========================================================================
\newpage
\section*{Appendix~A:\quad API Quick Reference}
\addcontentsline{toc}{section}{Appendix A: API Quick Reference}
% ========================================================================

The following table lists every \KAL\ function.

\bigskip
\noindent
\begin{longtable}{>{\tt}l l}
\hline
\multicolumn{1}{l}{\bf Function} &
\multicolumn{1}{l}{\bf Header} \\
\hline
\endfirsthead
\hline
\multicolumn{1}{l}{\bf Function} &
\multicolumn{1}{l}{\bf Header} \\
\hline
\endhead
\hline
\endfoot

\multicolumn{2}{l}{\bf Process Management} \\
\hline
kal\_getpid()            & kal\_process.h \\
kal\_process\_spawn()    & kal\_process.h \\
kal\_process\_kill()     & kal\_process.h \\
kal\_process\_exit()     & kal\_process.h \\
kal\_process\_reap()     & kal\_process.h \\
kal\_process\_wait()     & kal\_process.h \\
kal\_env\_get()          & kal\_process.h \\
kal\_env\_set()          & kal\_process.h \\
kal\_can\_exec()         & kal\_process.h \\
\hline

\multicolumn{2}{l}{\bf IPC / Sockets} \\
\hline
kal\_socket\_pair()      & kal\_ipc.h \\
kal\_close()             & kal\_ipc.h \\
kal\_write()             & kal\_ipc.h \\
kal\_read()              & kal\_ipc.h \\
kal\_poll()              & kal\_ipc.h \\
\hline

\multicolumn{2}{l}{\bf Memory} \\
\hline
kal\_malloc()            & kal\_memory.h \\
kal\_calloc()            & kal\_memory.h \\
kal\_realloc()           & kal\_memory.h \\
kal\_free()              & kal\_memory.h \\
\hline

\multicolumn{2}{l}{\bf Signals} \\
\hline
kal\_signal()            & kal\_signal.h \\
\hline

\multicolumn{2}{l}{\bf Threading} \\
\hline
kal\_mutex\_init()       & kal\_thread.h \\
kal\_mutex\_destroy()    & kal\_thread.h \\
kal\_mutex\_lock()       & kal\_thread.h \\
kal\_mutex\_trylock()    & kal\_thread.h \\
kal\_mutex\_unlock()     & kal\_thread.h \\
kal\_cond\_init()        & kal\_thread.h \\
kal\_cond\_destroy()     & kal\_thread.h \\
kal\_cond\_wait()        & kal\_thread.h \\
kal\_cond\_signal()      & kal\_thread.h \\
kal\_cond\_broadcast()   & kal\_thread.h \\
kal\_yield()             & kal\_thread.h \\
\hline

\multicolumn{2}{l}{\bf Terminal} \\
\hline
kal\_terminal\_save()    & kal\_terminal.h \\
kal\_terminal\_restore() & kal\_terminal.h \\
kal\_terminal\_raw()     & kal\_terminal.h \\
\hline

\multicolumn{2}{l}{\bf Filesystem} \\
\hline
kal\_fopen()             & kal\_filesys.h \\
kal\_fclose()            & kal\_filesys.h \\
kal\_fread()             & kal\_filesys.h \\
kal\_fwrite()            & kal\_filesys.h \\
kal\_fseek()             & kal\_filesys.h \\
kal\_open()              & kal\_filesys.h \\
kal\_fd\_write()         & kal\_filesys.h \\
kal\_fd\_close()         & kal\_filesys.h \\
kal\_mkstemp()           & kal\_filesys.h \\
kal\_mkdir()             & kal\_filesys.h \\
kal\_chmod()             & kal\_filesys.h \\
kal\_unlink()            & kal\_filesys.h \\
kal\_access()            & kal\_filesys.h \\
\hline

\multicolumn{2}{l}{\bf System} \\
\hline
kal\_get\_page\_size()          & kal\_system.h \\
kal\_get\_total\_memory()       & kal\_system.h \\
kal\_get\_available\_memory()   & kal\_system.h \\
kal\_get\_machine\_type()       & kal\_system.h \\
kal\_get\_os\_name()            & kal\_system.h \\
kal\_usleep()                   & kal\_system.h \\
\hline

\multicolumn{2}{l}{\bf Timing} \\
\hline
kal\_clock\_monotonic\_ns()  & kal\_time.h \\
kal\_clock\_monotonic()      & kal\_time.h \\
kal\_clock\_seconds()        & kal\_time.h \\
\hline

\end{longtable}

% ========================================================================
\section*{Appendix~B:\quad Configuration Limits}
\addcontentsline{toc}{section}{Appendix B: Configuration Limits}
% ========================================================================

The \KAL\ itself imposes no compile-time limits.  The following
limits apply to the \Mango\ kernel subsystems that use the \KAL:

\bigskip
\noindent
\begin{tabularx}{\textwidth}{lrlX}
\hline
\bf Symbol & \bf Value & & \bf Description \\
\hline
{\tt MACH\_PORT\_TABLE\_SIZE} & 256 & & Maximum Mach ports. \\
{\tt MACH\_PORT\_QUEUE\_MAX}  & 64  & & Messages per port queue. \\
{\tt TASK\_MAX}               & 128 & & Maximum tasks. \\
{\tt TASK\_PORT\_MAX}         & 64  & & Port rights per task. \\
{\tt THREAD\_MAX}             & 256 & & Threads per task. \\
{\tt IPC\_MAX\_SERVICES}      & 128 & & Named IPC services. \\
{\tt kal\_poll max}           & 64  & & Descriptors per poll call. \\
\hline
\end{tabularx}

\vfill

% -----------------------------------------------------------------------
%  Colophon
% -----------------------------------------------------------------------

\begin{center}
\small\itshape
This document was typeset in \TeX.\\
The body text is set in Computer Modern Roman,\\
and code examples in Computer Modern Typewriter.
\end{center}

\end{document}
